\documentclass[11pt]{article}

\usepackage[margin=1in]{geometry}
\usepackage{amsmath, amssymb, amsthm}
\usepackage{mathtools}
\usepackage{bbm}
\usepackage{hyperref}

\newtheorem{theorem}{Theorem}
\newtheorem{lemma}[theorem]{Lemma}
\newtheorem{proposition}[theorem]{Proposition}
\newtheorem{corollary}[theorem]{Corollary}
\newtheorem{definition}[theorem]{Definition}
\newtheorem{remark}[theorem]{Remark}
\newtheorem{assumption}[theorem]{Assumption}

\DeclareMathOperator*{\argmin}{arg\,min}
\DeclareMathOperator*{\argmax}{arg\,max}
\DeclareMathOperator{\softplus}{softplus}
\newcommand{\R}{\mathbb{R}}
\newcommand{\E}{\mathbb{E}}
\newcommand{\Lcal}{\mathcal{L}}
\newcommand{\Rcal}{\mathcal{R}}
\newcommand{\Hcal}{\mathcal{H}}
\newcommand{\Ncal}{\mathcal{N}}
\newcommand{\TV}{\mathrm{TV}}
\newcommand{\KL}{\mathrm{KL}}
\newcommand{\ind}{\mathbbm{1}}

\title{Theoretical Analysis of the PINN--Classical Router Loss Function: \\
Calibration, Bayes-Optimal Policy, and Domain Continuity}
\author{}
\date{}

\begin{document}
\maketitle

\begin{abstract}
We analyse the logistic loss function used to train a learned router that partitions
a PDE domain into regions solved by a physics-informed neural network (PINN) and a
classical numerical solver (e.g.\ finite-volume CFD, finite-element FEM, or
finite-difference FDM).  Since the PINN is frozen, the problem reduces to optimising
the router alone---a second-stage $\Rcal$-consistency question in the
predictor-rejector framework of Mao, Mohri, and Zhong~\cite{mao2024predictor}.  We
show that the logistic surrogate is \emph{classification-calibrated}: its pointwise
minimiser induces the Bayes-optimal binary assignment, a threshold rule on a
median-normalised residual--error cost.  We derive an $\Rcal$-consistency bound
relating target excess to surrogate excess, and analyse the domain continuity and
approximation properties of the resulting hybrid PINN--classical solution with
one-way Dirichlet coupling.  Our results are stated for a generic governing
operator $\Ncal$ on a discretised domain $\Omega \subset \R^d$ and apply
without modification to both the steady incompressible Navier--Stokes and Poisson
settings used in our experiments.
\end{abstract}

%% ========================================================================
\section{Problem Setup}
\label{sec:setup}
%% ========================================================================

\subsection{Domain and Solvers}

Let $\Omega \subset \R^d$ denote a bounded spatial domain discretised on a uniform
grid, with active region $\Omega_f = \{x \in \Omega : \mathrm{layout}(x) = 1\}$
(``active'' meaning the cells on which the PDE is to be solved---fluid cells for
Navier--Stokes, plate-minus-hole cells for Poisson, etc.) and $N \coloneqq |\Omega_f|
< \infty$.  Let $\Ncal$ denote the governing PDE operator (e.g.\ steady
incompressible Navier--Stokes, Poisson, Helmholtz, linear elasticity).  Two
solvers are available:

\begin{itemize}
\item A \textbf{PINN solver}~$h$, pre-trained and \emph{frozen}, producing an
  approximate solution field $\mathbf{u}_h \colon \Omega_f \to \R^m$ for each
  $x \in \Omega_f$, where $m$ is the number of solution components ($m = 3$ for
  steady NSE in 2D, with $\mathbf{u}_h = (u, v, p)$; $m = 1$ for scalar Poisson;
  $m = 2$ for plane-strain elasticity, etc.).  Its pointwise rejection cost
  combines the local PDE residual with an error-transport estimate (ETE) and is
  then median-normalised:
  \begin{equation}\label{eq:residual-def}
    R(x) \;\coloneqq\;
      \frac{\bigl\|\Ncal[\mathbf{u}_h](x)\bigr\| + \bigl|e(x)\bigr|}
           {\mathrm{median}_{x' \in \Omega_f}
            \bigl(\bigl\|\Ncal[\mathbf{u}_h](x')\bigr\| + \bigl|e(x')\bigr|\bigr)}
      \;\ge\; 0,
  \end{equation}
  where $e(x)$ is an estimate of the pointwise PINN error obtained by transporting
  the residual through a linearisation of the PDE (the error-transport equation
  $\Lcal_h\,e = \Ncal[\mathbf{u}_h]$, with $\Lcal_h$ the linearisation of $\Ncal$
  about $\mathbf{u}_h$).  Concrete ETE solvers depend on $\Ncal$: e.g.\ Jacobi
  sweeps of $\mathbf{u}_h \!\cdot\! \nabla e = \nu\,\Delta e$ for the
  advection--diffusion linearisation of NSE, an FFT pseudo-inverse of
  $-\Delta - k^2$ for Helmholtz, or a sparse direct solve of $-\Delta e$ for
  Poisson.  Median normalisation is robust to the heavy-tailed, right-skewed
  distribution typical of PDE error proxies: most of the domain has
  $R(x) \approx 1$, with a long tail where the PINN is inaccurate.  Combining the
  residual with the ETE $e$ corrects the well-known \emph{residual-to-error gap}:
  PDE residuals can be small at points where the PINN error is nonetheless large
  (and vice versa) because the differential operator can attenuate or amplify the
  underlying error~\cite{verfurth2013posteriori}.

\item A \textbf{classical solver}---a finite-volume / finite-element / finite-
  difference discretisation of $\Ncal$ (e.g.\ fractional-step CFD with
  Gauss--Seidel pressure Poisson for NSE, linear FEM for Poisson, direct
  factorisation for elasticity)---which solves the PDE to discretisation accuracy
  but at per-point computational cost~$\beta > 0$.  Since $R(x)$ has median~$1$,
  $\beta$ is directly interpretable independent of the PDE family: $\beta = 1.2$
  means ``invoke the classical solver wherever the residual--error cost exceeds
  $1.2\times$ the domain median.''
\end{itemize}

\subsection{Router}

A \textbf{router} is a function $s \colon \Omega_f \to \R$, parameterised by a
CNN~$s_\theta$, producing \emph{unbounded logits}:
\[
  s(x) \gg 0 \;\Longleftrightarrow\; \text{use the classical solver at } x,
  \qquad
  s(x) \ll 0 \;\Longleftrightarrow\; \text{use PINN at } x.
\]
The hypothesis class $\Rcal$ is the set of functions realisable by the architecture.
The hard assignment is obtained by thresholding: $\hat{d}(x) = \ind_{s(x) > 0}$.

\subsection{Target and Surrogate Losses}

The \emph{target (hard-assignment) loss} for a binary router
$d \colon \Omega_f \to \{0,1\}$:
\begin{equation}\label{eq:target}
  L(d) \;\coloneqq\; \frac{1}{N}\sum_{x \in \Omega_f}
    \bigl[\beta\,d(x) + (1 - d(x))\,R(x)\bigr].
\end{equation}
The \emph{surrogate (logistic) loss} for a continuous router
$s \colon \Omega_f \to \R$:
\begin{equation}\label{eq:surrogate}
  \boxed{\;
  \Lcal(s) \;\coloneqq\;
    \underbrace{\frac{1}{N}\sum_{x}
      \Bigl[\beta\,\varphi(s(x), 0)
        \;+\; R(x)\,\varphi(s(x), 1)\Bigr]}_{\text{logistic routing cost}}
    \;+\; \underbrace{\lambda_{\mathrm{tv}}\,\TV(s)}_{\text{spatial smoothness}},
  \;}
\end{equation}
where the logistic cost functions are
\begin{equation}\label{eq:phi}
  \varphi(s, 0) \;\coloneqq\; \ln(1 + e^{s}) = \softplus(s),
  \qquad
  \varphi(s, 1) \;\coloneqq\; \ln(1 + e^{-s}) = \softplus(-s),
\end{equation}
and $\TV(s) \coloneqq \frac{1}{N}\sum_{\text{adj.}~(x,x')} |s(x) - s(x')|$ is the
anisotropic total variation.  We write
$\bar{\Lcal}(s) \coloneqq
  \Lcal(s)\big|_{\lambda_{\mathrm{tv}}=0}$ for the
unregularised surrogate.

\begin{remark}[Logistic cost as a smooth relaxation]\label{rem:logistic}
The functions $\varphi(s,0)$ and $\varphi(s,1)$ are smooth, strictly positive, and
satisfy the asymptotic limits
\[
  \varphi(s,0) \to
    \begin{cases} s & s \to +\infty, \\ 0 & s \to -\infty, \end{cases}
  \qquad
  \varphi(s,1) \to
    \begin{cases} 0 & s \to +\infty, \\ -s & s \to -\infty. \end{cases}
\]
Thus $\varphi(s,j)$ behaves like the indicator cost $\ind_{s > 0}$ (for $j=0$) or
$\ind_{s \le 0}$ (for $j=1$) in the large-$|s|$ limit, but provides smooth gradients
for all finite~$s$.
\end{remark}

\begin{remark}[Connection to the predictor-rejector framework]\label{rem:framework}
In~\cite{mao2024predictor}, the abstention loss has the form
$\ell_{\mathrm{abs}}(h, r, x, y) = \ind_{h(x) \ne y}\,\ind_{r(x) > 0}
+ c\,\ind_{r(x) \le 0}$, where $h$ is a predictor and $r$ is a rejector.  In our
setting the PINN plays the role of~$h$ (first stage) and the router plays the role
of~$r$ (second stage), with the discrete classification error $\ind_{h(x) \ne y}$
replaced by the continuous residual--error cost~$R(x)$ and the abstention cost~$c$
by~$\beta$.  Crucially, since $h$ is \emph{frozen}, we do not optimise over any
predictor class~$\Hcal$.  The problem reduces to a pure \emph{$\Rcal$-consistency}
question for the router, with no joint $(\Hcal, \Rcal)$-optimisation required.
Nothing in the analysis below depends on the specific form of~$\Ncal$; the only
PDE-specific input to the proofs is the (nonnegative) field~$R(x)$.
\end{remark}

%% ========================================================================
\section{Calibration and Bayes-Optimal Policy}
\label{sec:calibration}
%% ========================================================================

Since the PINN is frozen, we only need to establish that minimising the logistic
surrogate over the router class~$\Rcal$ yields a good hard router.  We first analyse
the unregularised loss~$\bar{\Lcal}$ (i.e.\ $\lambda_{\mathrm{tv}} = 0$), which
decomposes into independent pointwise problems.

\subsection{Pointwise Minimiser and Calibration}

\begin{proposition}[Logistic surrogate is calibrated]\label{prop:calibrated}
For each $x \in \Omega_f$, the pointwise logistic cost
\begin{equation}\label{eq:ell-logistic}
  \ell(s; x) \;\coloneqq\;
    \beta\,\softplus(s) + R(x)\,\softplus(-s)
\end{equation}
is strictly convex in $s \in \R$ with unique minimiser
\begin{equation}\label{eq:s-star}
  s^*(x) \;=\; \ln\!\frac{R(x)}{\beta}.
\end{equation}
The induced binary decision $\ind_{s^*(x) > 0} = \ind_{R(x) > \beta}$ is the
Bayes-optimal hard assignment.
\end{proposition}

\begin{proof}
Differentiating~\eqref{eq:ell-logistic}:
\begin{equation}\label{eq:ell-deriv}
  \frac{\partial \ell}{\partial s}
  = \beta\,\sigma(s) - R(x)\,\sigma(-s)
  = \beta\,\sigma(s) - R(x)\,(1 - \sigma(s))
  = (\beta + R(x))\,\sigma(s) - R(x),
\end{equation}
where $\sigma(s) = 1/(1 + e^{-s})$ is the logistic sigmoid.  Setting to zero:
\begin{equation}
  \sigma(s^*) = \frac{R(x)}{\beta + R(x)},
  \qquad\text{hence}\qquad
  s^* = \ln\frac{R(x)}{\beta}.
\end{equation}
The second derivative is
\begin{equation}
  \frac{\partial^2 \ell}{\partial s^2}
  = (\beta + R(x))\,\sigma(s)\,(1 - \sigma(s)) > 0
  \qquad\text{for all } s \in \R,
\end{equation}
confirming strict convexity.  Finally,
$s^*(x) > 0 \iff R(x) > \beta$, so the binary decision
$\ind_{s^* > 0}$ coincides with the Bayes-optimal assignment.
\end{proof}

\begin{remark}[Value at the optimum]\label{rem:opt-value}
The minimum logistic cost at each $x$ is
\begin{equation}\label{eq:ell-star}
  \ell(s^*; x) = \beta\,\ln\!\frac{\beta + R}{\beta}
    + R\,\ln\!\frac{\beta + R}{R}
  = (\beta + R)\,H_2\!\!\left(\frac{\beta}{\beta + R}\right),
\end{equation}
where $H_2(p) = -p\ln p - (1-p)\ln(1-p)$ is the binary entropy in nats.  This is
strictly greater than the target minimum $\min\{\beta,\,R(x)\}$ for $\beta \ne R(x)$,
reflecting the smoothness cost of the logistic relaxation.  The gap vanishes as
$R(x)/\beta \to 0$ or $R(x)/\beta \to \infty$ (confident assignments).
\end{remark}

\subsection{Bayes-Optimal Router}

\begin{theorem}[Bayes-optimal router]\label{thm:bayes}
The unique minimiser of $\bar{\Lcal}$ over all $s \colon \Omega_f \to \R$ is
$s^*(x) = \ln(R(x)/\beta)$ pointwise.  The induced hard router
\begin{equation}\label{eq:bayes-router}
  d^*(x) =
  \begin{cases}
    0 & \text{if } R(x) < \beta \quad (\text{use PINN}), \\[3pt]
    1 & \text{if } R(x) > \beta \quad (\text{use CFD}),
  \end{cases}
\end{equation}
achieves optimal target value
$L(d^*) = \frac{1}{N}\sum_x \min\{\beta,\, R(x)\}$.
\end{theorem}

\begin{proof}
Since $\bar{\Lcal}(s) = \frac{1}{N}\sum_x \ell(s(x); x)$ decomposes as a sum of
independent pointwise terms, the global minimiser is the product of pointwise
minimisers from Proposition~\ref{prop:calibrated}.  The target value follows from
$\min\{\beta,\,R\} = \beta\,\ind_{R > \beta} + R\,\ind_{R \le \beta}$.
\end{proof}

\begin{remark}[Shape of the optimal router]\label{rem:shape}
The Bayes-optimal router is a \emph{level-set detector} on the residual--error
field.  The optimal classical-solver region is the superlevel set
\begin{equation}\label{eq:optimal-cfd}
  \Omega_C^* \;=\; \bigl\{x \in \Omega_f : R(x) > \beta\bigr\},
\end{equation}
and the optimal PINN region is
$\Omega_P^* = \{x \in \Omega_f : R(x) \le \beta\}$.

\textbf{Interpretation as a Bayes classifier.}
At each spatial point~$x$, the router performs a binary decision between two
``actions'': use PINN (cost $R(x)$) or use the classical solver (cost $\beta$).
The Bayes-optimal decision rule selects the action with lower cost---use the
classical solver if and only if the PINN residual--error cost exceeds the
classical-solver cost threshold.  This is the continuous-cost analogue of
Chow's optimal reject rule~\cite{chow1970optimum}, which abstains when posterior
uncertainty exceeds a cost~$c$.

\textbf{Physical interpretation.}
Since $R(x)$ is median-normalised, most of the domain has $R(x) \approx 1$.  The
long right tail---in regions where the PINN struggles, e.g.\ boundary layers and
recirculation zones for Navier--Stokes, sharp sources or geometric features for
Poisson and Helmholtz, or stress concentrations for elasticity---has $R(x) \gg 1$.
Both the residual term and the ETE term contribute to this tail: the residual
flags local equation imbalance, while the ETE captures regions where small
residuals nonetheless produce large transported error.  The router assigns these
challenging regions to the classical solver whenever $R(x) > \beta$.  The
parameter~$\beta$ controls the cost--accuracy trade-off: smaller~$\beta$ enlarges
$\Omega_C^*$ (more classical solving, more accurate, more expensive); larger
$\beta$ shrinks $\Omega_C^*$ (more PINN, cheaper, but potentially less accurate).
\end{remark}

\subsection{\texorpdfstring{$\Rcal$}{R}-Consistency Bound}

The non-trivial step is showing that a router with small \emph{surrogate} excess
(logistic loss close to optimal) also has small \emph{target} excess (binary cost
close to optimal).  Unlike the affine surrogate, the logistic surrogate is strictly
convex, so the consistency bound is of square-root type---standard for smooth
surrogates~\cite{bartlett2006convexity}.

\begin{lemma}[Pointwise surrogate gap]\label{lem:gap}
For each $x \in \Omega_f$, if the binary decision at~$s$ disagrees with the
Bayes-optimal decision (i.e.\ $\ind_{s > 0} \ne d^*(x)$), then
\begin{equation}\label{eq:gap}
  \ell(s; x) - \ell(s^*; x)
  \;\ge\;
  \frac{(\beta - R(x))^2}{2(\beta + R(x))}.
\end{equation}
\end{lemma}

\begin{proof}
If $\ind_{s > 0} \ne d^*(x)$, then $s$ and $s^* = \ln(R/\beta)$ have opposite signs.
Since $\ell(\cdot; x)$ is strictly convex with minimum at~$s^*$, and $s$ lies on the
opposite side of~$0$ from~$s^*$, we have $\ell(s; x) \ge \ell(0; x)$.  Therefore:
\begin{equation}
  \ell(s; x) - \ell(s^*; x)
  \;\ge\; \ell(0; x) - \ell(s^*; x)
  \;=\; (\beta + R)\,\ln 2 - (\beta + R)\,H_2\!\!\left(\frac{\beta}{\beta + R}\right).
\end{equation}
Writing $p = \beta/(\beta + R)$, the RHS equals
$(\beta + R)\,D_{\KL}(p \,\|\, \tfrac{1}{2})$,
where $D_{\KL}$ is the Kullback--Leibler divergence.  By Pinsker's
inequality~\cite{csiszar1967information}:
\begin{equation}
  D_{\KL}\!\left(p \,\Big\|\, \tfrac{1}{2}\right)
  \;\ge\; 2\left(p - \tfrac{1}{2}\right)^2
  = \frac{(\beta - R)^2}{2(\beta + R)^2}.
\end{equation}
Multiplying by $(\beta + R)$ yields~\eqref{eq:gap}.
\end{proof}

\begin{theorem}[$\Rcal$-consistency bound]\label{thm:r-consistency}
For any $s \colon \Omega_f \to \R$, define $d(x) = \ind_{s(x) > 0}$.  Then:
\begin{equation}\label{eq:r-consistency}
  L(d) - L(d^*)
  \;\le\;
  \sqrt{2(\beta + R_{\max})\,
    \bigl(\bar{\Lcal}(s) - \bar{\Lcal}(s^*)\bigr)},
\end{equation}
where $R_{\max} = \max_{x \in \Omega_f} R(x)$.
\end{theorem}

\begin{proof}
Let $S = \{x \in \Omega_f : d(x) \ne d^*(x)\}$ denote the set of
misclassified points, and $\delta(x) \coloneqq \ell(s(x); x) - \ell(s^*(x); x)$
the pointwise surrogate excess.

\textbf{Step 1.}  By Lemma~\ref{lem:gap}, for each $x \in S$:
\begin{equation}
  |\beta - R(x)|^2
  \;\le\; 2(\beta + R(x))\,\delta(x)
  \;\le\; 2(\beta + R_{\max})\,\delta(x).
\end{equation}

\textbf{Step 2.}  The target excess decomposes as:
\begin{equation}
  N\bigl(L(d) - L(d^*)\bigr) = \sum_{x \in S} |\beta - R(x)|.
\end{equation}
By the Cauchy--Schwarz inequality:
\begin{equation}
  \Bigl(\sum_{x \in S} |\beta - R(x)|\Bigr)^2
  \;\le\; |S| \cdot \sum_{x \in S} (\beta - R(x))^2
  \;\le\; |S| \cdot 2(\beta + R_{\max}) \sum_{x \in S} \delta(x).
\end{equation}

\textbf{Step 3.}  Since $|S| \le N$ and
$\sum_{x \in S} \delta(x) \le \sum_{x \in \Omega_f} \delta(x)
= N\,\bigl(\bar{\Lcal}(s) - \bar{\Lcal}(s^*)\bigr)$:
\begin{equation}
  N^2\bigl(L(d) - L(d^*)\bigr)^2
  \;\le\; N \cdot 2(\beta + R_{\max}) \cdot N\,
    \bigl(\bar{\Lcal}(s) - \bar{\Lcal}(s^*)\bigr).
\end{equation}
Dividing both sides by~$N^2$ and taking the square root
yields~\eqref{eq:r-consistency}.
\end{proof}

\begin{remark}[Interpretation of the bound]\label{rem:bound-interpretation}
The square-root form $\Delta L \le \sqrt{C \cdot \Delta\bar{\Lcal}}$ is
characteristic of smooth (strictly convex) surrogates---compare with the linear
bound $\Delta L \le C \cdot \Delta\bar{\Lcal}$ achievable for piecewise-linear
surrogates~\cite{bartlett2006convexity}.  The constant $\beta + R_{\max}$ reflects the
heavy-tailed residual distribution: extreme outliers (large $R_{\max}$) loosen the
bound but also produce large surrogate excess~$\delta(x)$ at those points,
partially compensating.

In practice, the bound is most relevant near the decision boundary $R(x) \approx
\beta$, where the cost of misclassification is small ($|\beta - R| \ll 1$) and the
square-root rate is not a limitation.  Far from the boundary, both the surrogate and
target costs strongly favour the correct decision.
\end{remark}

%% ========================================================================
\section{Domain Continuity and the Hybrid Solution}
\label{sec:continuity}
%% ========================================================================

We now analyse what happens \emph{after} the router produces a hard partition: the
domain is split into a classical-solver subdomain and a PINN subdomain, and a
hybrid solution is constructed by running the classical solver on its subdomain
with the PINN solution supplied as Dirichlet boundary data at the interface.

\subsection{Domain Decomposition}

\begin{definition}[Router-induced partition]\label{def:partition}
Given a binary router $d \colon \Omega_f \to \{0,1\}$, define:
\begin{align}
  \Omega_C &\;\coloneqq\; \{x \in \Omega_f : d(x) = 1\}
    && \text{(classical-solver region)}, \label{eq:omega-c}\\
  \Omega_P &\;\coloneqq\; \{x \in \Omega_f : d(x) = 0\}
    && \text{(PINN region)}, \label{eq:omega-p}\\
  \Gamma &\;\coloneqq\; \{x \in \Omega_C :
    \exists\, x' \in \Omega_P \text{ adjacent to } x\}
    && \text{(interface)}. \label{eq:gamma}
\end{align}
The interface $\Gamma \subset \Omega_C$ is the one-cell-wide layer of classical-
solver cells adjacent to the PINN region, identified via binary dilation
of~$\Omega_P$.
\end{definition}

\subsection{Hybrid Solution via One-Way Dirichlet Coupling}

\begin{definition}[Hybrid solution]\label{def:hybrid}
The hybrid PINN--classical solution
$\mathbf{u}_{\mathrm{hyb}} \colon \Omega_f \to \R^m$ is defined by:
\begin{equation}\label{eq:hybrid}
  \mathbf{u}_{\mathrm{hyb}}(x) \;=\;
  \begin{cases}
    \mathbf{u}_h(x) & \text{if } x \in \Omega_P
      \quad\text{(PINN solution)}, \\[3pt]
    \mathbf{u}_h(x) & \text{if } x \in \Gamma
      \quad\text{(Dirichlet coupling)}, \\[3pt]
    \mathbf{u}_C(x) & \text{if } x \in \Omega_C \setminus \Gamma
      \quad\text{(classical solution)},
  \end{cases}
\end{equation}
where $\mathbf{u}_C$ solves the discretised PDE $\Ncal[\mathbf{u}_C] = 0$ on
$\Omega_C$ with Dirichlet boundary conditions:
\begin{equation}\label{eq:cfd-bc}
  \mathbf{u}_C\big|_\Gamma = \mathbf{u}_h\big|_\Gamma
  \qquad\text{(PINN provides boundary data at the interface)}.
\end{equation}
All $m$ solution components use this coupling.  For Navier--Stokes ($m = 3$), this
means interface cells are reset to $p_h(x)$ during each pressure Poisson iteration
and to $(u_h, v_h)(x)$ during each momentum solve; for scalar Poisson ($m = 1$),
interface cells of the FEM solve are simply pinned to $u_h(x)$; analogous resets
apply to other PDE families.
\end{definition}

\begin{remark}[One-way vs.\ iterative coupling]\label{rem:coupling}
The coupling is \emph{one-directional}: the PINN supplies Dirichlet data to the
classical solver, but receives no feedback.  This differs from classical
overlapping or non-overlapping Schwarz domain
decomposition~\cite{toselli2005domain}, which iterates:
\begin{equation}\label{eq:schwarz-comparison}
  \underbrace{\text{PINN} \;\longrightarrow\; \text{classical}}_{\text{our approach: single sweep}}
  \qquad\text{vs.}\qquad
  \underbrace{\text{subdomain}_1 \;\rightleftharpoons\; \text{subdomain}_2}_{\text{Schwarz: iterate to convergence}}.
\end{equation}
Our single-sweep approach is justified when the PINN is sufficiently accurate at
the interface---a condition that the router loss directly controls (see
Theorem~\ref{thm:error-bound} and Remark~\ref{rem:loss-justification}).
\end{remark}

\subsection{Interface Continuity}

\begin{proposition}[Discrete continuity]\label{prop:continuity}
The hybrid solution $\mathbf{u}_{\mathrm{hyb}}$ has no jump discontinuity across
$\Gamma$: for every pair of adjacent cells $x \in \Gamma$ and $x' \in \Omega_P$,
\begin{equation}\label{eq:no-jump}
  \mathbf{u}_{\mathrm{hyb}}(x) = \mathbf{u}_h(x), \qquad
  \mathbf{u}_{\mathrm{hyb}}(x') = \mathbf{u}_h(x').
\end{equation}
\end{proposition}

\begin{proof}
By Definition~\ref{def:hybrid}: at $x \in \Gamma$,
$\mathbf{u}_{\mathrm{hyb}}(x) = \mathbf{u}_h(x)$; at $x' \in \Omega_P$,
$\mathbf{u}_{\mathrm{hyb}}(x') = \mathbf{u}_h(x')$.  Both values are evaluations of
the same (smooth) neural network $\mathbf{u}_h$ on adjacent grid cells, so no jump
occurs.
\end{proof}

\begin{remark}[Normal derivative discontinuity]\label{rem:normal-deriv}
While the solution values are continuous across~$\Gamma$, the \emph{normal derivatives}
are generally discontinuous: the gradient in $\Omega_P$ is determined by the PINN, while
the gradient in $\Omega_C$ is determined by the CFD solver.  This is characteristic of
non-overlapping domain decomposition with Dirichlet-only coupling (as opposed to
Robin coupling, which also matches fluxes).
\end{remark}

\subsection{Error Analysis of the Hybrid Solution}

Let $\mathbf{u}^*$ denote the exact PDE solution on~$\Omega$.

\begin{assumption}[Classical-solver stability]\label{ass:stability}
The classical solver on $\Omega_C$ satisfies a Dirichlet-data stability estimate:
for any two sets of boundary data $\mathbf{g}_1, \mathbf{g}_2$ on~$\Gamma$, the
corresponding solutions $\mathbf{u}_C^{(1)}, \mathbf{u}_C^{(2)}$ of the discretised
PDE on $\Omega_C$ satisfy
\begin{equation}\label{eq:stability}
  \bigl\|\mathbf{u}_C^{(1)} - \mathbf{u}_C^{(2)}\bigr\|_{\Omega_C \setminus \Gamma}
  \;\le\; C_s\,\bigl\|\mathbf{g}_1 - \mathbf{g}_2\bigr\|_\Gamma,
\end{equation}
where $C_s > 0$ is the stability constant and $\|\cdot\|_S$ denotes a suitable
norm (e.g.\ discrete $\ell^2$) over subset~$S$.
\end{assumption}

\begin{remark}
Assumption~\ref{ass:stability} is standard for well-posed PDEs.  For pure
elliptic problems (Poisson, scalar Helmholtz away from resonance, linear
elasticity), discrete maximum-principle and Lax--Milgram arguments give $C_s$
bounded by a constant depending only on coercivity and the inf--sup constant.
For the Stokes equations (the low-Reynolds linearisation of Navier--Stokes),
elliptic maximum-principle arguments give $C_s \le 1$.  For the full
Navier--Stokes at moderate Reynolds number (e.g.\ $\mathrm{Re} \approx 100$ in
our cylinder-flow experiments), $C_s$ is bounded and depends on the flow regime,
domain geometry, and mesh quality.
\end{remark}

\begin{theorem}[Hybrid solution error bound]\label{thm:error-bound}
Under Assumption~\ref{ass:stability}, the hybrid solution error decomposes as:
\begin{equation}\label{eq:error-bound}
  \bigl\|\mathbf{u}_{\mathrm{hyb}} - \mathbf{u}^*\bigr\|_{\Omega_f}
  \;\le\;
  \underbrace{\bigl\|\mathbf{u}_h - \mathbf{u}^*\bigr\|_{\Omega_P \cup \Gamma}
    \vphantom{\bigl\|}}_{\text{(I): PINN error}}
  \;+\; C_s\,\underbrace{\bigl\|\mathbf{u}_h - \mathbf{u}^*\bigr\|_{\Gamma}
    \vphantom{\bigl\|}}_{\text{(II): interface error}}
  \;+\; \underbrace{\epsilon_{\mathrm{disc}}
    \vphantom{\bigl\|}}_{\text{(III): classical discr.\ error}}.
\end{equation}
\end{theorem}

\begin{proof}
Decompose $\Omega_f$ into three disjoint subsets:
$\Omega_P$, $\Gamma$, and $\Omega_C \setminus \Gamma$.  By the triangle inequality:
\begin{equation}\label{eq:triangle}
  \bigl\|\mathbf{u}_{\mathrm{hyb}} - \mathbf{u}^*\bigr\|_{\Omega_f}
  \;\le\; \bigl\|\mathbf{u}_{\mathrm{hyb}} - \mathbf{u}^*\bigr\|_{\Omega_P \cup \Gamma}
  \;+\; \bigl\|\mathbf{u}_{\mathrm{hyb}} - \mathbf{u}^*\bigr\|_{\Omega_C \setminus \Gamma}.
\end{equation}

\textbf{Term~(I): PINN region and interface.}
On $\Omega_P \cup \Gamma$, $\mathbf{u}_{\mathrm{hyb}} = \mathbf{u}_h$ by
Definition~\ref{def:hybrid}, so
\begin{equation}\label{eq:term-I}
  \bigl\|\mathbf{u}_{\mathrm{hyb}} - \mathbf{u}^*\bigr\|_{\Omega_P \cup \Gamma}
  = \bigl\|\mathbf{u}_h - \mathbf{u}^*\bigr\|_{\Omega_P \cup \Gamma}.
\end{equation}

\textbf{Terms~(II) + (III): classical-solver interior.}
On $\Omega_C \setminus \Gamma$, $\mathbf{u}_{\mathrm{hyb}} = \mathbf{u}_C$, the
classical solution with boundary data $\mathbf{u}_h\big|_\Gamma$.  The exact
solution~$\mathbf{u}^*$, restricted to $\Omega_C$, solves the same PDE with
boundary data $\mathbf{u}^*\big|_\Gamma$ (and shared physical BCs on
$\partial\Omega_C \cap \partial\Omega$).  By uniqueness of PDE solutions on
$\Omega_C$, $\mathbf{u}^*\big|_{\Omega_C}$ is the classical solution with exact
boundary data.

Applying the stability estimate~\eqref{eq:stability} with
$\mathbf{g}_1 = \mathbf{u}_h\big|_\Gamma$ and
$\mathbf{g}_2 = \mathbf{u}^*\big|_\Gamma$, and adding the discretisation error:
\begin{equation}\label{eq:term-II-III}
  \bigl\|\mathbf{u}_C - \mathbf{u}^*\bigr\|_{\Omega_C \setminus \Gamma}
  \;\le\; C_s\,\bigl\|\mathbf{u}_h - \mathbf{u}^*\bigr\|_\Gamma
  \;+\; \epsilon_{\mathrm{disc}}.
\end{equation}

Substituting~\eqref{eq:term-I} and~\eqref{eq:term-II-III}
into~\eqref{eq:triangle} yields~\eqref{eq:error-bound}.
\end{proof}

\begin{remark}[Connection to the loss function]\label{rem:loss-justification}
Theorem~\ref{thm:error-bound} reveals the three error sources that the router loss
function~\eqref{eq:surrogate} controls:

\begin{enumerate}
\item[\textbf{(I)}] \textbf{PINN error in $\Omega_P$.}  The term
  $R(x)\,\varphi(s(x), 1)$ penalises high rejection cost in the PINN region
  (where $s(x) < 0$, so $\varphi(s,1) \approx |s|$ is large).  The cost
  $R(x) \propto \|\Ncal[\mathbf{u}_h](x)\| + |e(x)|$ combines the local PDE
  residual with the ETE estimate of the PINN error: both are proxies for
  $\|\mathbf{u}_h(x) - \mathbf{u}^*(x)\|$ via standard a~posteriori error
  estimates~\cite{verfurth2013posteriori}, and their sum tightens the proxy by
  flagging both \emph{equation imbalance} (residual) and \emph{transported error}
  (ETE).  By penalising $R(x)\,\varphi(s,1)$, the optimiser pushes high-cost
  points to $\Omega_C$, keeping only low-error points in~$\Omega_P$.

\item[\textbf{(II)}] \textbf{Interface PINN error (amplified by $C_s$).}  The PINN
  error at $\Gamma$ propagates into $\Omega_C$ via the Dirichlet coupling,
  amplified by~$C_s$.  Since the router loss penalises the residual--error cost
  across the entire PINN region (including points near the interface), the
  optimiser naturally places $\Gamma$ in low-cost areas, controlling this term.

\item[\textbf{(III)}] \textbf{Classical-solver discretisation error.}  The cost term
  $\beta\,\varphi(s(x), 0)$ budgets the computational expense of the classical
  solver.  The parameter~$\beta$ controls how much of~$\Omega_f$ is allocated to
  it.
\end{enumerate}
\end{remark}

\begin{corollary}[Compact form]\label{cor:compact}
Since $\Gamma \subset \Omega_P \cup \Gamma$, we can simplify to
\begin{equation}\label{eq:compact-bound}
  \bigl\|\mathbf{u}_{\mathrm{hyb}} - \mathbf{u}^*\bigr\|_{\Omega_f}
  \;\le\; (1 + C_s)\,\bigl\|\mathbf{u}_h - \mathbf{u}^*\bigr\|_{\Omega_P \cup \Gamma}
  \;+\; \epsilon_{\mathrm{disc}}.
\end{equation}
If the PINN is exact ($\mathbf{u}_h = \mathbf{u}^*$), then
$\mathbf{u}_{\mathrm{hyb}} = \mathbf{u}^*$ up to the classical-solver
discretisation error, regardless of the router partition.
\end{corollary}

%% ========================================================================
\section{Total Variation Regularisation}
\label{sec:tv}
%% ========================================================================

The training loss includes a total variation term
$\Lcal(s) = \bar{\Lcal}(s) + \lambda_{\mathrm{tv}}\,\TV(s)$
that promotes spatially contiguous classical-solver and PINN regions.  Since TV
acts on the raw logits~$s \in \R$ (not bounded to $[0,1]$), it penalises sharp
spatial gradients in the router output, promoting smooth transitions.

\begin{proposition}[TV-regularised consistency]\label{prop:tv}
Let $s^*_\lambda \in \argmin_{s \colon \Omega_f \to \R}
  \bigl[\bar{\Lcal}(s) + \lambda_{\mathrm{tv}}\,\TV(s)\bigr]$
and $d_\lambda(x) = \ind_{s^*_\lambda(x) > 0}$.  Then:
\begin{equation}\label{eq:tv-bound}
  L(d_\lambda) - L(d^*)
  \;\le\; \sqrt{2(\beta + R_{\max})\,\lambda_{\mathrm{tv}}\,\TV(s^*)},
\end{equation}
where $\TV(s^*)$ is the total variation of the unregularised optimal logits
$s^*(x) = \ln(R(x)/\beta)$.  In particular, the bias vanishes as
$\lambda_{\mathrm{tv}} \to 0$.
\end{proposition}

\begin{proof}
\textbf{Step 1.}  Since $s^*_\lambda$ minimises $\bar{\Lcal} + \lambda_{\mathrm{tv}}\,\TV$
and $s^*$ is feasible:
\begin{equation}\label{eq:tv-opt}
  \bar{\Lcal}(s^*_\lambda) + \lambda_{\mathrm{tv}}\,\TV(s^*_\lambda)
  \;\le\; \bar{\Lcal}(s^*) + \lambda_{\mathrm{tv}}\,\TV(s^*).
\end{equation}

\textbf{Step 2.}  Since $s^*$ minimises $\bar{\Lcal}$ globally:
$\bar{\Lcal}(s^*_\lambda) - \bar{\Lcal}(s^*) \ge 0$.

\textbf{Step 3.}  Rearranging~\eqref{eq:tv-opt} and using $\TV(s^*_\lambda) \ge 0$:
\begin{equation}\label{eq:surr-gap}
  0 \;\le\; \bar{\Lcal}(s^*_\lambda) - \bar{\Lcal}(s^*)
  \;\le\; \lambda_{\mathrm{tv}}\,\bigl(\TV(s^*) - \TV(s^*_\lambda)\bigr)
  \;\le\; \lambda_{\mathrm{tv}}\,\TV(s^*).
\end{equation}

\textbf{Step 4.}  Applying Theorem~\ref{thm:r-consistency}:
\begin{equation}\label{eq:tv-final}
  L(d_\lambda) - L(d^*)
  \;\le\; \sqrt{2(\beta + R_{\max})\,
    \bigl(\bar{\Lcal}(s^*_\lambda) - \bar{\Lcal}(s^*)\bigr)}
  \;\le\; \sqrt{2(\beta + R_{\max})\,\lambda_{\mathrm{tv}}\,\TV(s^*)}. \qedhere
\end{equation}
\end{proof}

\begin{remark}[TV and interface quality]\label{rem:tv-interface}
Beyond the abstract bias bound, TV regularisation provides a concrete benefit to the
domain continuity analysis of Section~\ref{sec:continuity}:
\begin{enumerate}
\item[(i)] \textbf{Smooth interfaces.}
  By penalising $|s(x) - s(x')|$ for adjacent cells, TV discourages jagged
  classical--PINN boundaries.  Smoother interfaces yield better-conditioned
  classical-solver problems---the stability constant~$C_s$ in
  Assumption~\ref{ass:stability} is smaller for geometrically regular subdomains
  than for highly fragmented ones.
\item[(ii)] \textbf{Contiguous subdomains.}
  TV promotes connected regions, which is numerically favourable: classical
  solvers (CFD/FEM/FDM) all perform best on compact, simply-connected subdomains
  where stencil or basis-function operations are well-defined everywhere.
\end{enumerate}
\end{remark}

\begin{remark}[Tension between TV and logistic sharpness]\label{rem:tension}
There is a fundamental tension between TV regularisation and the logistic loss.
The logistic loss rewards \emph{large} logits (confident decisions), since
$\softplus(s) \to 0$ as $s \to -\infty$ and $\softplus(-s) \to 0$ as $s \to +\infty$.
TV regularisation penalises \emph{spatial gradients} in~$s$, which must be large
near the classical--PINN boundary if the logits transition sharply.

In the coverage plot (loss vs.\ coverage), this manifests as a vertical gap between
the actual router operating point (diamond) and the theoretical optimum (star).
The diamond sits above the star because finite logits incur residual softplus cost
that the binary limit avoids.  The TV penalty prevents the router from sharpening its
logits to close this gap.  The parameter~$\lambda_{\mathrm{tv}}$ controls this
trade-off: larger values produce smoother but less confident routers.
\end{remark}

%% ========================================================================
\section{Summary of Guarantees}
\label{sec:summary}
%% ========================================================================

\begin{center}
\renewcommand{\arraystretch}{1.5}
\begin{tabular}{lll}
\hline
\textbf{Result} & \textbf{Statement} & \textbf{Ref.} \\
\hline
Calibration
  & $s^*(x) = \ln(R/\beta)$;\; $\ind_{s^* > 0} = \ind_{R > \beta}$ is Bayes-optimal
  & Prop.~\ref{prop:calibrated} \\
Bayes-optimal policy
  & $d^*(x) = \ind_{R(x) > \beta}$;\;
    $L^* = \frac{1}{N}\sum \min\{\beta, R\}$
  & Thm.~\ref{thm:bayes} \\
$\Rcal$-consistency
  & $\Delta L \le \sqrt{2(\beta + R_{\max})\,\Delta\bar{\Lcal}}$
  & Thm.~\ref{thm:r-consistency} \\
Interface continuity
  & $\mathbf{u}_{\mathrm{hyb}}$ is continuous across $\Gamma$
  & Prop.~\ref{prop:continuity} \\
Hybrid error bound
  & $\|e\|_\Omega \le (1{+}C_s)\,\|\mathbf{u}_h - \mathbf{u}^*\|_{\Omega_P \cup \Gamma}
    + \epsilon_{\mathrm{disc}}$
  & Cor.~\ref{cor:compact} \\
TV regularisation
  & bias $\le \sqrt{2(\beta + R_{\max})\,\lambda_{\mathrm{tv}}\,\TV(s^*)} \to 0$
  & Prop.~\ref{prop:tv} \\
\hline
\end{tabular}
\end{center}

\medskip

These results establish that the router loss function~\eqref{eq:surrogate} is
\emph{theoretically principled}:
\begin{enumerate}
\item \textbf{Calibrated}: the logistic surrogate is strictly convex at each point,
  with a unique minimiser $s^*(x) = \ln(R(x)/\beta)$ whose sign determines the
  Bayes-optimal binary assignment.  No additional thresholding heuristic is needed
  beyond $\mathrm{sign}(s)$.
\item \textbf{$\Rcal$-consistent}: the target excess (binary cost gap) is bounded by
  $\sqrt{C \cdot \Delta\bar{\Lcal}}$, where $C = 2(\beta + R_{\max})$.  This
  square-root rate is optimal for smooth surrogates~\cite{bartlett2006convexity} and
  vanishes as the surrogate excess tends to zero.
\item \textbf{Domain-continuous}: the hybrid PINN--classical solution constructed
  from the router partition is continuous at the interface by the Dirichlet
  coupling.  The hybrid error is controlled by the PINN error in $\Omega_P$
  (which the loss directly penalises), the interface error (amplified by the
  classical-solver stability constant $C_s$), and the classical discretisation
  error.
\item \textbf{Regularisation-safe}: TV regularisation introduces bounded bias
  $O(\sqrt{\lambda_{\mathrm{tv}}})$ that vanishes in the limit, while promoting
  smooth interfaces that improve classical-solver stability.  The tension between
  TV smoothness and logistic sharpness is controlled by~$\lambda_{\mathrm{tv}}$.
\item \textbf{PDE-agnostic}: every step of the analysis depends only on the
  nonnegative scalar field $R(x)$ and the abstract stability assumption on the
  classical solver.  The same theorems therefore apply to steady incompressible
  Navier--Stokes (with $\mathbf{u}_h = (u, v, p)$ and CFD as the classical
  solver), to Poisson and Helmholtz problems (scalar $\mathbf{u}_h$ with FEM or
  FDM), to linear elasticity, and to other steady PDEs admitting a Dirichlet
  stability bound.
\end{enumerate}

\bibliographystyle{plain}
\begin{thebibliography}{10}

\bibitem{mao2024predictor}
A.~Mao, M.~Mohri, and Y.~Zhong.
\newblock Predictor-rejector multi-class abstention: Theoretical analysis and
  algorithms.
\newblock In \emph{Proceedings of the 35th International Conference on
  Algorithmic Learning Theory (ALT)}, volume 237 of \emph{Proceedings of
  Machine Learning Research}, pages 1--47, 2024.

\bibitem{chow1970optimum}
C.~Chow.
\newblock On optimum recognition error and reject tradeoff.
\newblock \emph{IEEE Transactions on Information Theory}, 16(1):41--46, 1970.

\bibitem{bartlett2006convexity}
P.~L.~Bartlett, M.~I.~Jordan, and J.~D.~McAuliffe.
\newblock Convexity, classification, and risk bounds.
\newblock \emph{Journal of the American Statistical Association},
  101(473):138--156, 2006.

\bibitem{csiszar1967information}
I.~Csisz{\'a}r.
\newblock Information-type measures of difference of probability distributions
  and indirect observations.
\newblock \emph{Studia Scientiarum Mathematicarum Hungarica}, 2:299--318, 1967.

\bibitem{toselli2005domain}
A.~Toselli and O.~Widlund.
\newblock \emph{Domain Decomposition Methods --- Algorithms and Theory}.
\newblock Springer Series in Computational Mathematics, vol.~34, 2005.

\bibitem{verfurth2013posteriori}
R.~Verf\"urth.
\newblock \emph{A Posteriori Error Estimation Techniques for Finite Element
  Methods}.
\newblock Oxford University Press, 2013.

\end{thebibliography}

\end{document}
